% title:          Unique shift in a derivative equation
% area:           derivatives
% methods:        mean-value-theorem, construction
% difficulty:     medium
% difficulty-ai:  easy
% status:
% review:         human
% --- history: all optional, leave EMPTY if unknown ---
% origin:         imc
% origin-date:    2023-08-03
% origin-number:  Day 2, Problem 7
% also-in:
% references:     Earliest known appearance (to the best of our knowledge): IMC 2023, Second Day, August 3, 2023, Problem 7 (answer 1/(e-1)) - https://www.imc-math.org.uk/imc2023/imc2023-day2-questions.pdf
% history:        ai
\begin{problem}
    Let \( V := \left\{ f \in D^1\left(\mathbb{R},\mathbb{R}\right) \,\middle|\, f\left(0\right) = 0
    \wedge f\left(1\right) = 1 \right\} \). Compute
    \[
        S := \left\{ \alpha \in \mathbb{R} \,\middle|\, \forall f \in V \ \exists c \in
        \left]0;1\right[ \ f\left(c\right)+\alpha = f'\left(c\right) \right\}.
    \]
\end{problem}
\begin{solution}
    The idea is that \( f\left(c\right)+\alpha = f'\left(c\right) \) says that the derivative of
    \( \exp\left(-\_\right)\left(f+\alpha\right) \) vanishes at \( c \), while the mean value
    theorem on \( \left[0;1\right] \) gives a point where this derivative equals a number that
    depends only on \( \alpha \), since all the functions of \( V \) have the same values at \( 0 \)
    and at \( 1 \)!
    \\\\
    Let \( \alpha \in \mathbb{R} \) and \( f \in V \), and define
    \[
        g_{f,\alpha} \colon \mathbb{R} \longrightarrow \mathbb{R}, \quad x \longmapsto
        \exp\left(-x\right)\left(f\left(x\right)+\alpha\right),
    \]
    which is in \( D^1\left(\mathbb{R},\mathbb{R}\right) \), as a product of differentiable
    functions, and, for all \( x \in \mathbb{R} \),
    \[
        \left(g_{f,\alpha}\right)'\left(x\right) =
        \exp\left(-x\right)\left(f'\left(x\right)-\left(f\left(x\right)+\alpha\right)\right).
    \]
    Then there exists \( c \in \left]0;1\right[ \) with
    \[
        0 = \left(g_{f,\alpha}\right)'\left(c\right) =
        \exp\left(-c\right)\left(f'\left(c\right)-\left(f\left(c\right)+\alpha\right)\right)
    \]
    if and only if there exists \( c \in \left]0;1\right[ \) with \( f\left(c\right)+\alpha =
    f'\left(c\right) \), since \( \exp \) does not vanish. As \( \alpha \) and \( f \) were
    arbitrary, this implies that
    \[
        S = \left\{ \alpha \in \mathbb{R} \,\middle|\, \forall f \in V \ \exists c \in
        \left]0;1\right[ \ \left(g_{f,\alpha}\right)'\left(c\right) = 0 \right\}.
    \]
    Notice that, for all \( \alpha \in \mathbb{R} \) and \( f \in V \), the mean value theorem
    applied to \( g_{f,\alpha} \in D^1\left(\mathbb{R},\mathbb{R}\right) \) on
    \( \left[0;1\right] \) implies that there exists \( c \in \left]0;1\right[ \) with
    \[
        \left(g_{f,\alpha}\right)'\left(c\right) =
        \frac{g_{f,\alpha}\left(1\right)-g_{f,\alpha}\left(0\right)}{1-0} =
        g_{f,\alpha}\left(1\right)-g_{f,\alpha}\left(0\right) =
        \frac{1}{e}\left(1+\alpha\right)-\alpha,
    \]
    where the last equality holds because \( f \in V \). Since
    \( \frac{1}{e}\left(1+\alpha\right)-\alpha = 0 \) if and only if \( \alpha = \frac{1}{e-1} \),
    we get that, for all \( f \in V \), there exists \( c \in \left]0;1\right[ \) with
    \( \left(g_{f,\frac{1}{e-1}}\right)'\left(c\right) = 0 \), i.e., \( \frac{1}{e-1} \in S \).
    \\\\
    We claim that \( S = \left\{\frac{1}{e-1}\right\} \); by the description of \( S \) seen above,
    it suffices to let \( \alpha \in \mathbb{R} \setminus \left\{\frac{1}{e-1}\right\} \) and to
    show that there exists \( f_\alpha \in V \) such that, for all \( c \in \left]0;1\right[ \),
    \[
        0 \neq \left(g_{f_\alpha,\alpha}\right)'\left(c\right).
    \]
    The idea is the following. First, for any function \( f \in V \), the function
    \( g_{f,\alpha} \) takes the values \( \alpha \) at \( 0 \) and \( \frac{1+\alpha}{e} \) at
    \( 1 \). So, to make its derivative never zero in between, the easiest way is to take, between
    these two values, an affine function \( x \longmapsto ax+b \) with \( a \neq 0 \). Luckily, the
    values \( \alpha \) and \( \frac{1+\alpha}{e} \) are not equal, since \( \alpha \neq
    \frac{1}{e-1} \) (by choice of \( \alpha \)). So, the two points \( \left(0,\alpha\right) \) and
    \( \left(1,\frac{1+\alpha}{e}\right) \) have different heights, and hence the affine function
    joining them is not constant (i.e., it has a non-zero slope)! We can take the line between these
    two points and then solve for \( f_\alpha \), i.e., solve
    \( \exp\left(-x\right)\left(f_\alpha\left(x\right)+\alpha\right) =
    \left(\frac{1+\alpha}{e}-\alpha\right)x+\alpha \) for \( f_\alpha\left(x\right) \). This gives
    us a function \( f_\alpha \in V \) such that \( g_{f_\alpha,\alpha} \) has a non-zero derivative
    in between. Solving for such an \( f_\alpha \) explains the following definition.
    \\\\
    Let
    \[
        f_\alpha \colon \mathbb{R} \longrightarrow \mathbb{R}, \quad x \longmapsto
        \left(\left(\frac{1+\alpha}{e}-\alpha\right)x+\alpha\right)\exp\left(x\right)-\alpha,
    \]
    which is in \( D^1\left(\mathbb{R},\mathbb{R}\right) \), as a product of differentiable
    functions minus a constant. Then \( f_\alpha\left(0\right) = \alpha-\alpha = 0 \) and
    \( f_\alpha\left(1\right) = \frac{1+\alpha}{e}\,e-\alpha = 1 \), so \( f_\alpha \in V \). Now,
    for all \( x \in \mathbb{R} \),
    \[
        \left(f_\alpha\right)'\left(x\right) =
        \left(\frac{1+\alpha}{e}-\alpha+\left(\frac{1+\alpha}{e}-\alpha\right)x+\alpha\right)\exp\left(x\right)
    \]
    \[
        = \left(\left(\frac{1+\alpha}{e}-\alpha\right)x+\frac{1+\alpha}{e}\right)\exp\left(x\right),
    \]
    so that
    \[
        \left(g_{f_\alpha,\alpha}\right)'\left(x\right) =
        \exp\left(-x\right)\left(\left(f_\alpha\right)'\left(x\right)-\left(f_\alpha\left(x\right)+\alpha\right)\right)
    \]
    \[
        =
        \exp\left(-x\right)\left(\left(\left(\frac{1+\alpha}{e}-\alpha\right)x+\frac{1+\alpha}{e}\right)\exp\left(x\right)-\left(\left(\frac{1+\alpha}{e}-\alpha\right)x+\alpha\right)\exp\left(x\right)\right)
    \]
    \[
        =
        \left(\frac{1+\alpha}{e}-\alpha\right)x+\frac{1+\alpha}{e}-\left(\left(\frac{1+\alpha}{e}-\alpha\right)x+\alpha\right)
        = \frac{1+\alpha}{e}-\alpha
    \]
    after cancelling the two terms \( \left(\frac{1+\alpha}{e}-\alpha\right)x \), i.e.,
    \( \left(g_{f_\alpha,\alpha}\right)'\left(x\right) = \frac{1+\alpha}{e}-\alpha \neq 0 \) for all
    \( x \in \mathbb{R} \), since \( \alpha \neq \frac{1}{e-1} \) (by choice of \( \alpha \)). In
    particular, \( \left(g_{f_\alpha,\alpha}\right)'\left(c\right) \neq 0 \) for all \( c \in
    \left]0;1\right[ \). Hence, \( S = \left\{\frac{1}{e-1}\right\} \).
\end{solution}
